% GENERATED FILE. Edit canonical lessons, then run npm run generate:books.
% canonical-source-hash: fnv1a64:086803a7dc51ed9c
% canonical-lessons: TE-C173-varanda, TE-C173-peradu, TE-C173-hotalu, TE-C173-pakkana, TE-C173-postaphisu

\chapter{Veranda, Backyard, Restaurant, Beside, Post office}
\label{ch:te-veranda-backyard-restaurant-beside-post-office}
\hlchaptermodality{173}

\begin{chapteropening}
\textbf{By the end of this chapter:} \emph{I can say a veranda, a backyard, a restaurant, beside and a post office.}

Complete the last lesson of chapter 173: I can say a veranda, a backyard, a restaurant, beside and a post office.
\end{chapteropening}

\section[\texorpdfstring{\te{వరండా} (varaṇḍā)}{వరండా (varaṇḍā)}]{\te{వరండా} (varaṇḍā) — a veranda}

\label{lesson:TE-C173-varanda}



\begin{quote}
Before the new one: say the Telugu for every year, then the Telugu for forever.
\end{quote}

\subsection*{You'll want to know: \te{వరండా}}
\textbf{\te{వరండా}} — \emph{varaṇḍā} — \textquotedblleft{}a veranda\textquotedblright{}.

\begin{grammarlens}[title={What you've built}]
One more word for this chapter.
\end{grammarlens}

\subsection*{Your turn}
\begin{itemize}
\raggedright
  \item \emph{Say it:} \emph{varaṇḍā}
  \item \emph{Say it:} \emph{varaṇḍā}, once more
  \item \emph{Say it:} say \emph{varaṇḍā}
  \item \emph{Recall:} say the Telugu for every year, then the Telugu for forever, then say \emph{varaṇḍā} again
\end{itemize}

\begin{tcolorbox}[breakable,colback=teal!4,colframe=teal!35!black,title={Before you move on}]
What does \te{వరండా} mean? (\textquotedblleft{}A veranda\textquotedblright{}.) Say it once more.
\end{tcolorbox}
\section[\texorpdfstring{\te{పెరడు} (peraḍu)}{పెరడు (peraḍu)}]{\te{పెరడు} (peraḍu) — a backyard}

\label{lesson:TE-C173-peradu}



\begin{quote}
Before the new one: say the Telugu for forever, then the Telugu for a veranda.
\end{quote}

\subsection*{You'll want to know: \te{పెరడు}}
\textbf{\te{పెరడు}} — \emph{peraḍu} — \textquotedblleft{}a backyard\textquotedblright{}.

\begin{grammarlens}[title={What you've built}]
One more word for this chapter.
\end{grammarlens}

\subsection*{Your turn}
\begin{itemize}
\raggedright
  \item \emph{Say it:} \emph{peraḍu}
  \item \emph{Say it:} \emph{peraḍu}, once more
  \item \emph{Say it:} say \emph{peraḍu}
  \item \emph{Recall:} say the Telugu for forever, then the Telugu for a veranda, then say \emph{peraḍu} again
\end{itemize}

\begin{tcolorbox}[breakable,colback=teal!4,colframe=teal!35!black,title={Before you move on}]
What does \te{పెరడు} mean? (\textquotedblleft{}A backyard\textquotedblright{}.) Say it once more.
\end{tcolorbox}
\section[\texorpdfstring{\te{హోటలు} (hōṭalu)}{హోటలు (hōṭalu)}]{\te{హోటలు} (hōṭalu) — a restaurant, an eating house}

\label{lesson:TE-C173-hotalu}



\begin{quote}
Before the new one: say the Telugu for a veranda, then the Telugu for a backyard.
\end{quote}

\subsection*{You'll want to know: \te{హోటలు}}
\textbf{\te{హోటలు}} — \emph{hōṭalu} — \textquotedblleft{}a restaurant, an eating house\textquotedblright{}.

\begin{grammarlens}[title={What you've built}]
One more word for this chapter.
\end{grammarlens}

\subsection*{Your turn}
\begin{itemize}
\raggedright
  \item \emph{Say it:} \emph{hōṭalu}
  \item \emph{Say it:} \emph{hōṭalu}, once more
  \item \emph{Say it:} say \emph{hōṭalu}
  \item \emph{Recall:} say the Telugu for a veranda, then the Telugu for a backyard, then say \emph{hōṭalu} again
\end{itemize}

\begin{tcolorbox}[breakable,colback=teal!4,colframe=teal!35!black,title={Before you move on}]
What does \te{హోటలు} mean? (\textquotedblleft{}A restaurant, an eating house\textquotedblright{}.) Say it once more.
\end{tcolorbox}
\section[\texorpdfstring{\te{పక్కన} (pakkana)}{పక్కన (pakkana)}]{\te{పక్కన} (pakkana) — beside, next to}

\label{lesson:TE-C173-pakkana}



\begin{quote}
Before the new one: say the Telugu for a backyard, then the Telugu for a restaurant.
\end{quote}

\subsection*{You'll want to know: \te{పక్కన}}
\textbf{\te{పక్కన}} — \emph{pakkana} — \textquotedblleft{}beside, next to\textquotedblright{}.

\begin{grammarlens}[title={What you've built}]
One more word for this chapter.
\end{grammarlens}

\subsection*{Your turn}
\begin{itemize}
\raggedright
  \item \emph{Say it:} \emph{pakkana}
  \item \emph{Say it:} \emph{pakkana}, once more
  \item \emph{Say it:} say \emph{pakkana}
  \item \emph{Recall:} say the Telugu for a backyard, then the Telugu for a restaurant, then say \emph{pakkana} again
\end{itemize}

\begin{tcolorbox}[breakable,colback=teal!4,colframe=teal!35!black,title={Before you move on}]
What does \te{పక్కన} mean? (\textquotedblleft{}Beside, next to\textquotedblright{}.) Say it once more.
\end{tcolorbox}
\section[\texorpdfstring{\te{పోస్టాఫీసు} (pōsṭāphīsu)}{పోస్టాఫీసు (pōsṭāphīsu)}]{\te{పోస్టాఫీసు} (pōsṭāphīsu) — a post office}

\label{lesson:TE-C173-postaphisu}



\begin{quote}
Before the new one: say the Telugu for a restaurant, then the Telugu for beside.
\end{quote}

\subsection*{You'll want to know: \te{పోస్టాఫీసు}}
\textbf{\te{పోస్టాఫీసు}} — \emph{pōsṭāphīsu} — \textquotedblleft{}a post office\textquotedblright{}.

\begin{grammarlens}[title={What you've built}]
One more word for this chapter.
\end{grammarlens}

\subsection*{Your turn}
\begin{itemize}
\raggedright
  \item \emph{Say it:} \emph{pōsṭāphīsu}
  \item \emph{Say it:} \emph{pōsṭāphīsu}, once more
  \item \emph{Say it:} say \emph{pōsṭāphīsu}
  \item \emph{Recall:} say the Telugu for a restaurant, then the Telugu for beside, then say \emph{pōsṭāphīsu} again
\end{itemize}

\begin{tcolorbox}[breakable,colback=teal!4,colframe=teal!35!black,title={Before you move on}]
What does \te{పోస్టాఫీసు} mean? (\textquotedblleft{}A post office\textquotedblright{}.) Say it once more.
\end{tcolorbox}
